\documentclass[10pt,a4paper]{amsart}
\usepackage[margin=19mm]{geometry}
\usepackage{amsmath,amssymb,mathtools}
\usepackage[scr=rsfs,cal=euler]{mathalfa}
\usepackage{xfrac}
\usepackage[unicode,hidelinks,bookmarks=false]{hyperref}
\newcommand{\ie}{i.e.\ }
\newcommand{\cf}{cf.\ }
\newcommand{\resp}{respectively\ }
\newcommand{\Subsection}{\subsection}
\providecommand{\nobreakdash}{\nobreak\hbox{-}}
\setlength{\emergencystretch}{2em}
\multlinegap0pt
\usepackage{amssymb}
\usepackage{mathtools}
\newtheorem{theorem}{Theorem}
\newcommand{\C}{\mathbb C}
\title{A Note on Additive Diameter}
\author{Ernesto Ingrosso}
\date{Source: arXiv:2609.06023v1 (CC BY 4.0)}
\begin{document}
\maketitle

\noindent\emph{Source and licence.} A Note on Additive Diameter. Version arXiv:2609.06023v1 (CC BY 4.0), distributed by arXiv under the Creative Commons Attribution 4.0 International licence, http://creativecommons.org/licenses/by/4.0/, as recorded in the arXiv licence metadata for this exact version and shown on the abstract page https://arxiv.org/abs/2609.06023v1. Full licence notice: \texttt{licenses/arxiv-2609.06023-CC-BY-4.0.txt}.\par\smallskip

\noindent\textit{Editorial scope.} Counted unit: the article's theorem thm:rightineq (source0.tex lines 530--540). The article's complete original proof of it is delivered with it (source0.tex lines 542--607, one \texttt{proof} environment of 66 lines). No mathematics is written, repaired or re-derived: the statement and the proof are the article's own lines, in the article's own notation, and no retained line is substituted or re-flowed.  No further source-local context is needed: every cross-reference of every retained line resolves inside the two delivered spans.  Nothing was omitted from any retained span, so the package carries no omission record.  No retained line contains a citation, so the wrapper carries no bibliography and no bibliography entry.  No retained line contains a figure environment, an image-inclusion command or drawing code, so no image is delivered and the package declares no asset dependency.  Editorial formatting only, in addition: an AMS \texttt{amsart} preamble that restates the article's own declarations and macros used by the retained lines, namely the environments \texttt{theorem} on separate counters, together with the standard packages those lines need.  The retained environments print in this document as Theorem 1 in order of appearance, whereas the article prints the same items with the numbering of its own full document.  The complete native source of the article as distributed in the arXiv e-print for this exact version is bundled unchanged as \texttt{sources/209484/source0.tex}.
\begin{theorem}\label{thm:rightineq}
Let $G$ be a complex algebraic group, let
$\rho:G\to\operatorname{GL}(V)$ be a finite-dimensional representation,
and let $\mathfrak g=\operatorname{Lie}(G)$. Then, for every subspace
$U\leq V$,
\[
 \operatorname{diam}^{G}_{+}(V,U)
 \leq
 \operatorname{diam}^{\mathfrak g}_{+}(V,U).
\]
\end{theorem}

\begin{proof}
Assume that $\operatorname{diam}^{\mathfrak g}_{+}(V,U)=d+1<\infty.$
Choose $x_1,\ldots,x_d\in\mathfrak g$, and put $X_i=d\rho(x_i),$
so that $V=U+X_1U+\cdots+X_dU.$

We may therefore choose elements $u_1,\ldots,u_r\in U$
and $v_1,\ldots,v_s\in U$, where $i_1,\ldots,i_s\in\{1,\ldots,d\},$
such that $u_1,\ldots,u_r,
 X_{i_1}v_1,\ldots,X_{i_s}v_s$
is a basis of $V$. For $t\in\C$ put $w_j(t)
 =
 \bigl(\operatorname{exp}(tX_{i_j})-I\bigr)v_j.$
 
Since
\[
 \operatorname{exp}(tX_{i_j})-I
 =
 tX_{i_j}+t^2\frac{X_{i_j}^2}{2!}
 +t^3\frac{X_{i_j}^3}{3!}+\cdots,
\]
we have $w_j(t)
 =
 tX_{i_j}v_j+t^2z_j(t)$
for some analytic map $z_j:\mathbb C\to V$.

Fixing a basis of $V$ and taking determinants, multilinearity gives
\[
\begin{aligned}
 &\det\bigl(
 u_1,\ldots,u_r,
 w_1(t),\ldots,w_s(t)
 \bigr)
\\
 &\qquad =
 t^s
 \det\bigl(
 u_1,\ldots,u_r,
 X_{i_1}v_1,\ldots,X_{i_s}v_s
 \bigr)
 +t^{s+1}h(t)
\end{aligned}
\]
for some analytic function $h(t)$.

The first determinant is nonzero. Hence, for some $t\neq0$, $u_1,\ldots,u_r,
 w_1(t),\ldots,w_s(t)$
is again a basis of $V$.

But $u_j\in U$ and $w_j(t)
 =
 \operatorname{exp}(tX_{i_j})v_j-v_j
 \in
 U+\operatorname{exp}(tX_{i_j})U,$ and so $V
 =
 U+\operatorname{exp}(tX_1)U+\cdots+
 \operatorname{exp}(tX_d)U.$

Now set $g_i=\operatorname{exp}(tx_i)\in G.$
Since $\rho(g_i)=\operatorname{exp}(tX_i),$ we obtain $V
 =
 U+\rho(g_1)U+\cdots+\rho(g_d)U.$
Thus $\operatorname{diam}^{G}_{+}(V,U)
 \leq d+1
 =
 \operatorname{diam}^{\mathfrak g}_{+}(V,U).$
\end{proof}

\end{document}
